% GENERATED FILE. Edit canonical lessons, then run npm run generate:books.
% canonical-source-hash: fnv1a64:170b00e558335c32
% canonical-lessons: GU-C134-bahadur, GU-C134-sant, GU-C134-khus, GU-C134-udas, GU-C134-bhukhyu

\chapter{Brave, Calm, Happy, Sad, Hungry}
\label{ch:gu-brave-calm-happy-sad-hungry}
\hlchaptermodality{134}

\begin{chapteropening}
\textbf{By the end of this chapter:} \emph{I can say brave, calm, happy, sad and hungry.}

Complete the last lesson of chapter 134: I can say brave, calm, happy, sad and hungry.
\end{chapteropening}

\section[\texorpdfstring{\gu{બહાદુર} (bahādur)}{બહાદુર (bahādur)}]{\gu{બહાદુર} (bahādur) — brave}

\label{lesson:GU-C134-bahadur}



\begin{quote}
Before the new one: say the Gujarati for clever, then the Gujarati for lazy.
\end{quote}

\subsection*{You'll want to know: \gu{બહાદુર}}
\textbf{\gu{બહાદુર}} (\emph{bahādur}) — \textquotedblleft{}brave\textquotedblright{}.

\begin{grammarlens}[title={What you've built}]
One more word for this chapter.
\end{grammarlens}

\subsection*{Your turn}
\begin{itemize}
\raggedright
  \item \emph{Say it:} \emph{bahādur}
  \item \emph{Say it:} \emph{bahādur}, once more
  \item \emph{Say it:} say \emph{bahādur}
  \item \emph{Recall:} say the Gujarati for clever, then the Gujarati for lazy, then say \emph{bahādur} again
\end{itemize}

\begin{tcolorbox}[breakable,colback=teal!4,colframe=teal!35!black,title={Before you move on}]
What does \gu{બહાદુર} mean? (\textquotedblleft{}Brave\textquotedblright{}.) Say it once more.
\end{tcolorbox}
\section[\texorpdfstring{\gu{શાંત} (śānt)}{શાંત (śānt)}]{\gu{શાંત} (śānt) — calm, quiet}

\label{lesson:GU-C134-sant}



\begin{quote}
Before the new one: say the Gujarati for lazy, then the Gujarati for brave.
\end{quote}

\subsection*{You'll want to know: \gu{શાંત}}
\textbf{\gu{શાંત}} (\emph{śānt}) — \textquotedblleft{}calm, quiet\textquotedblright{}.

\begin{grammarlens}[title={What you've built}]
One more word for this chapter.
\end{grammarlens}

\subsection*{Your turn}
\begin{itemize}
\raggedright
  \item \emph{Say it:} \emph{śānt}
  \item \emph{Say it:} \emph{śānt}, once more
  \item \emph{Say it:} say \emph{śānt}
  \item \emph{Recall:} say the Gujarati for lazy, then the Gujarati for brave, then say \emph{śānt} again
\end{itemize}

\begin{tcolorbox}[breakable,colback=teal!4,colframe=teal!35!black,title={Before you move on}]
What does \gu{શાંત} mean? (\textquotedblleft{}Calm, quiet\textquotedblright{}.) Say it once more.
\end{tcolorbox}
\section[\texorpdfstring{\gu{ખુશ} (khuś)}{ખુશ (khuś)}]{\gu{ખુશ} (khuś) — happy}

\label{lesson:GU-C134-khus}



\begin{quote}
Before the new one: say the Gujarati for brave, then the Gujarati for calm.
\end{quote}

\subsection*{You'll want to know: \gu{ખુશ}}
\textbf{\gu{ખુશ}} (\emph{khuś}) — \textquotedblleft{}happy\textquotedblright{}.

\begin{grammarlens}[title={What you've built}]
One more word for this chapter.
\end{grammarlens}

\subsection*{Your turn}
\begin{itemize}
\raggedright
  \item \emph{Say it:} \emph{khuś}
  \item \emph{Say it:} \emph{khuś}, once more
  \item \emph{Say it:} say \emph{khuś}
  \item \emph{Recall:} say the Gujarati for brave, then the Gujarati for calm, then say \emph{khuś} again
\end{itemize}

\begin{tcolorbox}[breakable,colback=teal!4,colframe=teal!35!black,title={Before you move on}]
What does \gu{ખુશ} mean? (\textquotedblleft{}Happy\textquotedblright{}.) Say it once more.
\end{tcolorbox}
\section[\texorpdfstring{\gu{ઉદાસ} (udās)}{ઉદાસ (udās)}]{\gu{ઉદાસ} (udās) — sad}

\label{lesson:GU-C134-udas}



\begin{quote}
Before the new one: say the Gujarati for calm, then the Gujarati for happy.
\end{quote}

\subsection*{You'll want to know: \gu{ઉદાસ}}
\textbf{\gu{ઉદાસ}} (\emph{udās}) — \textquotedblleft{}sad\textquotedblright{}.

\begin{grammarlens}[title={What you've built}]
One more word for this chapter.
\end{grammarlens}

\subsection*{Your turn}
\begin{itemize}
\raggedright
  \item \emph{Say it:} \emph{udās}
  \item \emph{Say it:} \emph{udās}, once more
  \item \emph{Say it:} say \emph{udās}
  \item \emph{Recall:} say the Gujarati for calm, then the Gujarati for happy, then say \emph{udās} again
\end{itemize}

\begin{tcolorbox}[breakable,colback=teal!4,colframe=teal!35!black,title={Before you move on}]
What does \gu{ઉદાસ} mean? (\textquotedblleft{}Sad\textquotedblright{}.) Say it once more.
\end{tcolorbox}
\section[\texorpdfstring{\gu{ભૂખ્યું} (bhūkhyũ)}{ભૂખ્યું (bhūkhyũ)}]{\gu{ભૂખ્યું} (bhūkhyũ) — hungry}

\label{lesson:GU-C134-bhukhyu}



\begin{quote}
Before the new one: say the Gujarati for happy, then the Gujarati for sad.
\end{quote}

\subsection*{You'll want to know: \gu{ભૂખ્યું}}
\textbf{\gu{ભૂખ્યું}} (\emph{bhūkhyũ}) — \textquotedblleft{}hungry\textquotedblright{}.

\begin{grammarlens}[title={What you've built}]
One more word for this chapter.
\end{grammarlens}

\subsection*{Your turn}
\begin{itemize}
\raggedright
  \item \emph{Say it:} \emph{bhūkhyũ}
  \item \emph{Say it:} \emph{bhūkhyũ}, once more
  \item \emph{Say it:} say \emph{bhūkhyũ}
  \item \emph{Recall:} say the Gujarati for happy, then the Gujarati for sad, then say \emph{bhūkhyũ} again
\end{itemize}

\begin{tcolorbox}[breakable,colback=teal!4,colframe=teal!35!black,title={Before you move on}]
What does \gu{ભૂખ્યું} mean? (\textquotedblleft{}Hungry\textquotedblright{}.) Say it once more.
\end{tcolorbox}
